\SourceLeaf{395}%
\[
\begin{array}{l}
M\\
M-2(a-r\cos A)\sin(A''-A')\\
M-2(a'-r\cos A')\sin(A-A'')\\
M-2(a''-r\cos A'')\sin(A'-A)
\end{array}
\]
Man hat also offenbar eine Gleichung der Form
\[
\gamma+\delta r+\varepsilon rr+\zeta r^3=\theta\theta-uu=4\alpha\alpha\beta\beta
\]
mit den gegebenen Größen \(\gamma,\delta,\varepsilon,\zeta\). Aus der Bedingung für das Maximum ergibt sich dann offenbar \(r\) aus der folgenden quadratischen Gleichung
\[
\delta+2\varepsilon r+3\zeta rr=0
\]
\par
Noch leichter lassen sich die Koeffizienten dieser Gleichung mit folgender Überlegung finden. Bei jeder Ellipse, die die vier Seiten des Vierecks berührt und deren Mittelpunkt die Abszisse \(r\) hat, ist das Vierfache des Produkts aus den Quadraten der großen und der kleinen Halbachse allgemein
\[
=\gamma+\delta r+\varepsilon rr+\zeta r^3
\]
also muss dieser Ausdruck notwendig \(=0\) werden, wenn man für \(r\) einen Wert einsetzt, der einer der drei oben betrachteten verschwindenden Ellipsen entspricht. Diese drei Werte sind die Abstände der Mittelpunkte der beiden Diagonalen des Vierecks und des Mittelpunkts der Strecke, die die Schnittpunkte der beiden Seitenpaare des Vierecks verbindet, vom Ursprung der Abszissen. Ich bezeichne diese drei Punkte mit \(C,D,E\) und ihre Abszissen mit \(c,d,e\). Dann muss offenbar
\[
r^3+\tfrac{\varepsilon}{\zeta}rr+\tfrac{\delta}{\zeta}r+\tfrac{\gamma}{\zeta}
\]
mit dem Produkt \((r-c)(r-d)(r-e)\) identisch sein. Folglich lautet die obige quadratische Gleichung
\[
3rr-2r(c+d+e)+cd+ce+de=0
\]
Die Wurzeln sind
\[
\tfrac{c+d+e}{3}+\tfrac13\BareRoot{(cc+dd+ee+cd+ce+de)}
\]
